\documentclass{article}
\usepackage[UTF8]{ctex}
\usepackage{amsmath}
\usepackage{amssymb}
\usepackage{graphicx}
\usepackage{geometry}
\geometry{a4paper, margin=2cm}

\newcommand{\problemname}[1]{\begin{center}\Large\bfseries #1\end{center}}
\usepackage{enumitem}

\begin{document}
\problemname{输出奇偶性}

\subsection*{题目描述}
设计程序，判断输入的数字是奇数还是偶数。例如输入 16，那么需要输出 oushu；如果输入 19，那么要输出 jishu。

\subsection*{输入格式}
输入只有一行，包含一个整数 shu。它在 int 类型变量可以存储的范围内。

\subsection*{输出格式}
如果这个数是偶数（能被 2 整除），输出 oushu；否则输出 jishu。

\subsection*{样例输入}
\begin{verbatim}
16
\end{verbatim}

\subsection*{样例输出}
\begin{verbatim}
oushu
\end{verbatim}


\end{document}
